主方案采用 A 边界（年末 SOC 为 $1200\,\mathrm{kWh}$）、因果负荷预测、线性锚点光伏映射、相对日初承诺的一次性调整结算，以及 `M1\_M6` 策略。以 2025 年 1 月 1 日 $6000\,\mathrm{kWh}$ 为起点并全年连续运行时，主方案实际成本为 $16\,373\,508.75$ 元；不利用后续预报调整的 M0 成本为 $17\,109\,568.25$ 元。因此，三次预测更新和价值触发调整使年度成本减少 $736\,059.50$ 元，降幅为 $4.30\%$。

该比较说明后续预报的价值并非来自任意频繁修改计划，而是来自在可调整的剩余时段内，用新信息权衡调整费用、紧急购电风险与储能状态。`M1\_M6` 中的 M6 是对每个更新时点信息价值的可解释执行记录；在同一目标、约束和货币容差下，它不应被和 M1 解释成两项可相加的独立收益。

年末 SOC 改为 $6000\,\mathrm{kWh}$ 的边界敏感性不改变策略排序，但使主策略成本增加 $2\,262.88$ 元，说明主结论不依赖于以过低年末 SOC 虚减成本。正式 `result3.xlsx` 按题设只含 2--12 月，共 $334$ 日，其区间成本为 $14\,512\,748.53$ 元；该值服务于输出核对，不能误称为全年成本。
